% ch03_control.tex -- Control Structures
\chapter{Control Structures --- Conditions and Loops}
\label{ch:controle}

\begin{intuition}
Imagine a cooking recipe: ``If the batter is too runny, add flour;
otherwise, bake for 30 minutes.'' A program works in exactly the same
way: it makes \emph{decisions} and \emph{repeats} steps according to conditions.
Control structures are the heart of every scientific algorithm.
\end{intuition}

% ============================================================
\section{Conditions: \py{if}, \py{elif}, \py{else}}
\index{if}\index{elif}\index{else}\index{condition}

In science, we often classify according to thresholds. Consider the
\emph{phase transitions}\index{phase transition} of water:

\begin{codeblock}[title=\textbf{Python}]
\begin{minted}{python}
temperature = -5  # in degrees Celsius

if temperature < 0:
    etat = "solide"
elif temperature < 100:
    etat = "liquide"
else:
    etat = "gazeux"

print(f"At {temperature} C, water is {etat}.")
\end{minted}
\end{codeblock}

\begin{codeoutput}
\begin{verbatim}
At -5 C, water is solide.
\end{verbatim}
\end{codeoutput}

\begin{remark}
Python evaluates conditions \emph{in order}. As soon as a condition is true,
the corresponding block is executed and the remaining ones are skipped.
\end{remark}

% --- Flowchart TikZ ---
\begin{figure}[ht]
\centering
\begin{tikzpicture}[
    decision/.style={diamond, draw, fill=blue!10, text width=3cm,
        align=center, inner sep=2pt},
    block/.style={rectangle, draw, fill=green!10, rounded corners,
        minimum height=1cm, text width=3cm, align=center},
    arr/.style={->, thick, >=stealth}
]
\node[block] (start) {Start};
\node[decision, below=1cm of start] (d1) {$T < 0$?};
\node[block, right=2.5cm of d1] (sol) {State = solid};
\node[decision, below=1.5cm of d1] (d2) {$T < 100$?};
\node[block, right=2.5cm of d2] (liq) {State = liquid};
\node[block, below=1.5cm of d2] (gaz) {State = gas};

\draw[arr] (start) -- (d1);
\draw[arr] (d1) -- node[above]{yes} (sol);
\draw[arr] (d1) -- node[left]{no} (d2);
\draw[arr] (d2) -- node[above]{yes} (liq);
\draw[arr] (d2) -- node[left]{no} (gaz);
\end{tikzpicture}
\caption{Flowchart of an \py{if/elif/else} structure for phase transitions.}
\label{fig:flowchart-if}
\end{figure}

\subsection{Comparison operators and logical operators}
\index{comparison operator}\index{logical operator}

\begin{definition}[Comparison operators]
The operators \py{==}, \py{!=}, \py{<}, \py{>}, \py{<=}, \py{>=} return
a Boolean (\py{True} or \py{False}). Conditions can be combined with
\py{and}, \py{or}, \py{not}.
\end{definition}

\begin{codeblock}[title=\textbf{Python}]
\begin{minted}{python}
pH = 3.2

if pH < 7 and pH >= 0:
    print("Acidic solution")
elif pH == 7:
    print("Neutral solution")
elif pH > 7 and pH <= 14:
    print("Basic solution")
else:
    print("Invalid pH value")
\end{minted}
\end{codeblock}

\begin{codeoutput}
\begin{verbatim}
Acidic solution
\end{verbatim}
\end{codeoutput}

% ============================================================
\section{The \py{while} loop}
\index{while}\index{loop}

The \py{while} loop repeats a block \emph{as long as a condition is true}.
In science, it is often used for \emph{convergence}\index{convergence} algorithms.

\begin{example}[Convergence of a sequence]
The sequence $u_{n+1} = \frac{1}{2}\left(u_n + \frac{a}{u_n}\right)$ converges to $\sqrt{a}$
(Heron's method)\index{Heron's method}.
\end{example}

\begin{codeblock}[title=\textbf{Python}]
\begin{minted}{python}
a = 2.0
u = 1.0       # initial value
tolerance = 1e-10
iterations = 0

while abs(u**2 - a) > tolerance:
    u = 0.5 * (u + a / u)
    iterations += 1

print(f"sqrt({a}) = {u}")
print(f"Convergence in {iterations} iterations")
\end{minted}
\end{codeblock}

\begin{codeoutput}
\begin{verbatim}
sqrt(2.0) = 1.4142135623730951
Convergence in 4 iterations
\end{verbatim}
\end{codeoutput}

\begin{warning}
A \py{while} loop whose condition \emph{never} becomes false runs
forever! Always add a safeguard:
\begin{minted}{python}
max_iter = 10000
while condition and iterations < max_iter:
    ...
\end{minted}
\end{warning}

% ============================================================
\section{The \py{for} loop and \py{range()}}
\index{for}\index{range}

The \py{for} loop iterates over an \emph{iterable} (list, string, \py{range}, etc.).

\begin{codeblock}[title=\textbf{Python}]
\begin{minted}{python}
# Compute the sum of squares from 1 to 10
somme = 0
for i in range(1, 11):
    somme += i**2

print(f"Sum of squares from 1 to 10 = {somme}")
# Verification: n(n+1)(2n+1)/6
print(f"Formula: {10 * 11 * 21 // 6}")
\end{minted}
\end{codeblock}

\begin{codeoutput}
\begin{verbatim}
Sum of squares from 1 to 10 = 385
Formula: 385
\end{verbatim}
\end{codeoutput}

\begin{definition}[\py{range(start, stop, step)}]
\py{range(a, b, s)} generates integers from \py{a} to \py{b-1} with a step of \py{s}.
Note: the upper bound \py{b} is \emph{excluded}.
\end{definition}

\begin{codeblock}[title=\textbf{Python}]
\begin{minted}{python}
# Display multiples of 3 between 0 and 15
for n in range(0, 16, 3):
    print(n, end=" ")
\end{minted}
\end{codeblock}

\begin{codeoutput}
\begin{verbatim}
0 3 6 9 12 15
\end{verbatim}
\end{codeoutput}

% ============================================================
\section{Nested loops}
\index{nested loop}

\begin{codeblock}[title=\textbf{Python}]
\begin{minted}{python}
# Multiplication table (excerpt)
for i in range(1, 4):
    for j in range(1, 4):
        print(f"{i} x {j} = {i*j:2d}", end="   ")
    print()  # newline
\end{minted}
\end{codeblock}

\begin{codeoutput}
\begin{verbatim}
1 x 1 =  1   1 x 2 =  2   1 x 3 =  3
2 x 1 =  2   2 x 2 =  4   2 x 3 =  6
3 x 1 =  3   3 x 2 =  6   3 x 3 =  9
\end{verbatim}
\end{codeoutput}

% ============================================================
\section{\py{break} and \py{continue}}
\index{break}\index{continue}

\begin{codeblock}[title=\textbf{Python}]
\begin{minted}{python}
# Find the first prime number > 50
n = 51
while True:
    est_premier = True
    for d in range(2, int(n**0.5) + 1):
        if n % d == 0:
            est_premier = False
            break       # no need to test other divisors
    if est_premier:
        print(f"First prime > 50: {n}")
        break           # exit the while loop
    n += 1
\end{minted}
\end{codeblock}

\begin{codeoutput}
\begin{verbatim}
First prime > 50: 53
\end{verbatim}
\end{codeoutput}

\begin{codeblock}[title=\textbf{Python}]
\begin{minted}{python}
# Skip negative values in measurements
mesures = [2.3, -1.0, 4.5, -0.3, 3.8, 7.1]
somme = 0
compteur = 0
for m in mesures:
    if m < 0:
        continue   # skip to the next iteration
    somme += m
    compteur += 1

print(f"Mean (positive values) = {somme/compteur:.2f}")
\end{minted}
\end{codeblock}

\begin{codeoutput}
\begin{verbatim}
Mean (positive values) = 4.43
\end{verbatim}
\end{codeoutput}

% ============================================================
\section{Common errors}

\begin{errmark}
\begin{enumerate}
    \item \textbf{Indentation error}: Python uses indentation to delimit
    blocks. Mixing spaces and tabs causes an \py{IndentationError}.
    Rule: always use \textbf{4 spaces}.

    \item \textbf{Infinite loop}: forgetting to increment the control variable
    in a \py{while} leads to a loop that never stops.
    \begin{minted}{python}
# ERROR: i never changes!
i = 0
while i < 10:
    print(i)    # prints 0 indefinitely
    \end{minted}

    \item \textbf{Off-by-one error}: \py{range(1, 10)} produces
    integers from 1 to \textbf{9}, not 10! To include 10, write \py{range(1, 11)}.

    \item \textbf{Using \py{=} instead of \py{==}}: \py{=} is assignment,
    \py{==} is comparison. Writing \py{if x = 5} causes a \py{SyntaxError}.
\end{enumerate}
\end{errmark}

% ============================================================
\section{Mini-project: Radioactive Decay}

\begin{projet}
The law of radioactive decay\index{radioactive decay} gives the number
of remaining atoms:
\[
    N(t) = N_0 \times 0.5^{\,t / t_{1/2}}
\]
where $N_0$ is the initial number of atoms and $t_{1/2}$ is the half-life.

\textbf{Objective}: simulate the decay of carbon-14 ($t_{1/2} = 5730$ years)
and display the number of remaining atoms every 1000 years.
\end{projet}

\begin{codeblock}[title=\textbf{Python}]
\begin{minted}{python}
N0 = 1_000_000       # initial number of atoms
demi_vie = 5730      # half-life of C-14 in years
seuil = N0 * 0.01    # stop when less than 1% remains

t = 0
N = N0
print(f"{'Time (yrs)':>12} {'N remaining':>12} {'% remaining':>10}")
print("-" * 36)

while N > seuil:
    pourcentage = N / N0 * 100
    print(f"{t:>12} {N:>12.0f} {pourcentage:>9.1f}%")
    t += 1000
    N = N0 * 0.5 ** (t / demi_vie)

print(f"\nAfter {t} years, less than 1% of the atoms remain.")
\end{minted}
\end{codeblock}

\begin{codeoutput}
\begin{verbatim}
 Time (yrs)   N remaining  % remaining
------------------------------------
           0      1000000     100.0%
        1000       886076      88.6%
        2000       785128      78.5%
        3000       695685      69.6%
        4000       616441      61.6%
        5000       546274      54.6%
        ...
       37000        12055       1.2%

After 38000 years, less than 1% of the atoms remain.
\end{verbatim}
\end{codeoutput}

% ============================================================
\section{Exercises}

\begin{exercise}[$\star$ --- BMI Classification]
Write a program that asks for a person's weight (kg) and height (m),
computes the BMI ($\text{BMI} = \text{weight} / \text{height}^2$) and displays the
category: underweight ($<18.5$), normal ($18.5$--$25$), overweight ($25$--$30$),
obese ($\geq 30$).
\end{exercise}

\begin{exercise}[$\star$ --- Countdown]
Use a \py{while} loop to display a countdown from 10 to 0,
then display ``Liftoff!''.
\end{exercise}

\begin{exercise}[$\star\star$ --- Fibonacci Sequence]
Compute the first 20 terms of the Fibonacci sequence\index{Fibonacci}
($F_0 = 0$, $F_1 = 1$, $F_{n} = F_{n-1} + F_{n-2}$) with a \py{for} loop.
Verify that the ratio $F_n / F_{n-1}$ converges to the golden ratio
$\varphi \approx 1.618$.
\end{exercise}

\begin{exercise}[$\star\star$ --- Sieve of Eratosthenes]
Use nested loops to find all prime numbers less than 100.
\end{exercise}

\begin{exercise}[$\star\star\star$ --- Random Walk Simulation]
Simulate a 1D random walk\index{random walk}: at each step, the particle moves
by $+1$ or $-1$ with equal probability. After $N = 1000$ steps,
display the final position. Repeat 100 times and compute the mean
distance $\langle |x| \rangle$. Compare with the theoretical prediction
$\sqrt{N} \approx 31.6$.
\emph{Hint:} use \py{import random} and \py{random.choice([-1, 1])}.
\end{exercise}

% ============================================================
\begin{keyformulas}{Chapter Summary: Control Structures}
\begin{itemize}
    \item \textbf{Conditions}: \py{if condition:} / \py{elif condition:} / \py{else:}
    \item \textbf{\py{while} loop}: repeats as long as the condition is true
    \item \textbf{\py{for} loop}: iterates over an iterable (\py{range()}, list, etc.)
    \item \textbf{\py{range(a, b, s)}}: integers from \py{a} to \py{b-1}, step \py{s}
    \item \textbf{\py{break}}: immediately exits the loop
    \item \textbf{\py{continue}}: skips to the next iteration
    \item \textbf{Indentation}: 4 spaces, always consistent
    \item \textbf{Safeguard}: always set a maximum number of iterations for \py{while}
\end{itemize}
\end{keyformulas}
